\chapter{Anaphylaxis}

\section*{Definition and Overview}
Anaphylaxis is a severe, potentially life-threatening allergic reaction that comes on quickly and affects more than one part of the body at once. It happens when the immune system overreacts to a trigger, such as a food, insect sting, or medicine, and releases large amounts of chemical mediators that cause symptoms within minutes to hours.

Common warning signs include hives, swelling of the face or throat, difficulty breathing, and a drop in blood pressure. Anaphylaxis is a medical emergency that requires immediate treatment with adrenaline and urgent medical care.

\section*{Epidemiology}
Anaphylaxis can occur at any age, but it is most common in children and young adults. Peanut and tree nut allergy is among the most frequent causes in children, while in adults medicines and insect stings account for a larger share of reactions.

Rates of hospital admission for anaphylaxis have risen over recent decades in many countries, although most reactions do not prove fatal. People with asthma, especially when poorly controlled, and those who have had anaphylaxis before are at higher risk of more severe reactions.

\section*{Aetiology and Pathophysiology}
\begin{itemize}
    \item \textbf{Food triggers:} peanuts, tree nuts, milk, eggs, fish, shellfish, wheat, and sesame are the most common culprits
    \item \textbf{Insect stings:} bee and wasp stings are frequent causes in older children and adults
    \item \textbf{Medicines:} antibiotics (especially penicillins), non-steroidal anti-inflammatory drugs (NSAIDs), and some other medicines
    \item \textbf{Other triggers:} latex, some anaesthetic drugs, and rare causes such as exercise-induced reactions or reactions to some blood products
    \item \textbf{Mechanism:} in most cases the immune system produces antibodies against the trigger; on next exposure these antibodies activate allergy cells, which rapidly release histamine and other chemicals, causing blood vessels to widen, fluid to leak into tissues, and airways to narrow
    \item \textbf{Non-immune reactions:} some reactions produce the same symptoms without a detectable immune antibody, but both types are treated in the same way
\end{itemize}

\section*{Clinical Features}
\begin{itemize}
    \item Hives (urticaria), itching, and flushed or pale skin
    \item Swelling of the lips, tongue, face, or throat
    \item Difficulty breathing, wheezing, a hoarse voice, or trouble speaking
    \item A feeling of faintness, dizziness, or a fast heart rate
    \item Nausea, vomiting, abdominal pain, or diarrhoea
    \item Collapse, loss of consciousness, or a strong sense that something is seriously wrong
\end{itemize}

\section*{Differential Diagnosis}
\begin{itemize}
    \item A severe asthma attack, with wheezing but without other allergic features
    \item A vasovagal (faint) episode, with feeling faint but without skin or throat symptoms
    \item A panic attack or anxiety, which can cause breathlessness and a racing heart
    \item Septic shock or a severe infection causing low blood pressure
    \item Food poisoning causing vomiting and abdominal pain
    \item Mild drug or food reactions that affect only one part of the body and are not anaphylaxis
    \item Angioedema caused by medicines such as ACE inhibitors, which swells tissues without a full allergic reaction
\end{itemize}

\section*{When to Seek Medical Care}
Anaphylaxis is a medical emergency. Do not wait to see whether symptoms settle, and do not delay treatment while seeking advice.

\textbf{Call 000 (or local emergency services) for:}
\begin{itemize}
    \item Difficulty breathing, wheezing, or swelling of the tongue or throat
    \item Faintness, collapse, or loss of consciousness
    \item Any suspected anaphylactic reaction, even if adrenaline has already been given
    \item A reaction that involves more than one body area, such as the skin and the stomach, or the skin and breathing, and is getting worse
\end{itemize}

\section*{Diagnosis and Investigations}
\begin{itemize}
    \item \textbf{Clinical assessment:} anaphylaxis is diagnosed from the rapid onset of symptoms and a likely trigger; there is no single test that proves it at the time of the reaction
    \item \textbf{Blood test:} a serum tryptase level taken soon after the reaction may support the diagnosis
    \item \textbf{Allergy testing:} skin prick tests or blood tests for specific antibodies (IgE) performed weeks later by an allergy specialist to identify the trigger
    \item \textbf{Challenge testing:} in selected cases, a controlled food or drug challenge in a supervised setting
    \item \textbf{Review by an allergist:} to confirm the cause and provide a written emergency action plan
\end{itemize}

\section*{Management}
\begin{itemize}
    \item \textbf{Emergency first aid:} lie the person flat with the legs raised, unless breathing is very difficult, when sitting up may help; call for an ambulance immediately
    \item \textbf{Adrenaline:} the first-line treatment, given as an autoinjector into the outer thigh; repeat after 5 minutes if symptoms do not improve
    \item \textbf{Hospital care:} oxygen, intravenous fluids, and other medicines such as antihistamines, corticosteroids, and inhaled bronchodilators as needed, with observation for several hours because reactions can return
    \item \textbf{Emergency action plan:} a written plan describing the triggers, symptoms, and how and when to use adrenaline
    \item \textbf{Ongoing management:} avoidance of known triggers, carrying an adrenaline autoinjector at all times, and regular review by an allergy specialist
\end{itemize}

\section*{Complications}
\begin{itemize}
    \item A biphasic reaction, in which symptoms return hours after the first episode, which is why hospital observation is advised
    \item Severe throat swelling or bronchospasm that blocks breathing
    \item Low blood pressure and shock, which can damage organs if not treated quickly
    \item Injury from a fall or accident during the reaction
    \item Psychological effects, including anxiety about the next reaction and reduced confidence in managing daily life
\end{itemize}

\section*{Prognosis and Outcomes}
With prompt adrenaline treatment, most people with anaphylaxis recover fully. Deaths are rare but occur most often when adrenaline is delayed or the person has severe asthma. About one in five people may experience a second wave of symptoms after the first wave settles, so observation in hospital is important. Most people who have had anaphylaxis go on to live active lives by avoiding their triggers, carrying adrenaline, and following a personalised action plan. Outcomes depend on recognising symptoms early and treating quickly.

\section*{Prevention and Self-Care}
\begin{itemize}
    \item Identify triggers with an allergy specialist and avoid them carefully, including reading food labels every time
    \item Carry an adrenaline autoinjector at all times and know how to use it
    \item Have a written anaphylaxis action plan and make sure family, school, or workplace contacts understand it
    \item Tell health professionals about your allergies before any medicine, anaesthetic, or vaccination
    \item If you are allergic to bee or wasp stings, avoid areas where insects gather and consider venom immunotherapy on specialist advice
    \item Manage asthma well, because poorly controlled asthma makes severe reactions more likely
\end{itemize}

\section*{Sources and Further Reading}
The following resources provide reliable, peer-reviewed and patient-orientated information on anaphylaxis.

\begin{itemize}
    \item NHS. \textit{Overview of anaphylaxis, its causes, and emergency treatment.} https://www.nhs.uk/conditions/
    \item Mayo Clinic. \textit{Guide to the symptoms and emergency management of anaphylaxis.} https://www.mayoclinic.org/
    \item World Health Organization. \textit{International information on allergy and anaphylaxis awareness.} https://www.who.int/
    \item MSD Manual. \textit{Professional and patient information on anaphylaxis.} https://www.msdmanuals.com/
    \item MedlinePlus. \textit{Health information on anaphylaxis.} https://medlineplus.gov/
\end{itemize}
